\subsection{導入：なぜ経済性が必要なのか}

ソフトウェアは多くの組織にとって、単なる業務補助ではなく、事業目標を達成するための手段である。売上を伸ばす、顧客体験を改善する、業務の再作業を減らす、保守の手戻りを減らす、競争力を高めるといった目標は、ソフトウェアの技術的な設計や運用判断と直接つながる。

工学経済は「お金の学問」だけではなく、「選択の学問」である。すなわち、限られた資源をこの技術的活動へ投入すべきか、それとも別の活動へ投入した方がよいかを考える。したがって、技術的に実現可能であることと、経済的に受け入れ可能であることのバランスを探す活動である。

\begin{pointbox}{重要}
技術的に優れた案が、組織にとって最良の案とは限らない。高性能な設計、最新の技術、完全な作り直しは、費用、納期、リスク、運用負荷、組織の無形資産との適合性を考えると不利になることがある。
\end{pointbox}

ソフトウェア工学経済は、開発前の投資判断だけでなく、ソフトウェア製品ライフサイクル全体に関係する。例として、次のような問いはすべて経済的な側面を持つ。
\begin{itemize}
  \item ある機能は買うべきか、作るべきか。
  \item ある要求を今回の範囲に含めるべきか。
  \item どの設計が最も費用対効果が高いか。
  \item クラウド環境で、応答時間を満たしつつ過剰な運用費を避ける構成はどれか。
  \item リスクに基づくテストはどこまで行えば十分か。
  \item 技術的負債の大きいコードを、リファクタリングするか、再開発するか、そのまま維持するか。
\end{itemize}

価値は金銭に限らない。企業市民としての責任、従業員の幸福、環境への配慮、顧客ロイヤルティ、ブランド信頼など、数値化しにくい価値も意思決定に影響する。このため、ソフトウェア工学経済では財務的な指標だけでなく、無形の基準も扱う必要がある。
